% Modernised reading — fonds Alexandre Grothendieck, shelfmark 134-8. Of the
% folder's 86 pages, only page 63 is transcribed here (batch-04.fr.tex, which
% covers pages 61-80 and carries that one leaf); this reading covers exactly
% that page.
%
% This is an INTERPRETATION, not a transcription. In any dispute the
% transcription governs: whoever cites Grothendieck must cite batch-04.fr.tex.
%
% Pass: Opus 5.5 (claude-opus-5-5), 2026-10-03 — first pass, unchecked against
% the pages by a human. It was made on Opus 5.5 and has not been measured
% against the reference reading of folder 115 (made on Opus 5). What is said of
% the circulation of Pursuing Stacks, and of the people named on the leaf, is
% cited from memory and was not checked against the sources.
%
% Scope (issue #46). Folder 134-8 belongs to Pursuing Stacks, of which
% G. Maltsiniotis has given the edition; the folder is left to that edition. The
% one leaf transcribed is the one the edition does not carry: typed, signed
% instructions for copying « Modelizing Story », pp. 259-381 of the typescript.
% There is no mathematics on it, so the reading is short and has no spine
% section; the four failure modes of the skill have nothing to bite on.
%
% Departures and choices, each footnoted:
%  - the leaf is undated; « été 1983 » is inferred (from the PS and from the
%    page numbers), not read;
%  - the leaf asks for pp. 259-381, while the inventory titles the folder by
%    chapter VII (pp. 555-593, Oct.-Nov. 1983): why it was filed here is not
%    known, and the reading does not guess;
%  - « Modelizing Story » is kept as his name for the part; no mapping onto
%    the chapters of the printed edition is asserted;
%  - his deletions and additions (photocopie -> xérographie, etc.) are folded
%    into the text and named in footnotes;
%  - « Gwynned » is reproduced as typed and footnoted (Gwynedd);
%  - the two private addresses withheld by the transcription stay withheld.
%    The university address of the first recipient, kept by the
%    transcription, is kept.
\documentclass[11pt,a4paper]{article}
\input{../preamble/grothendieck}

\folder{134-8}
\pages{61}{80}
\dating{[1983]}
\watermark{Édition de démonstration\\interprétation personnelle de l'œuvre}
\foldertitle{[Chapitre VII (pages 555 à 593) : Linearization of homotopy types] : tapuscrit annoté (22/10-12/11/[1983]), notes manuscrites (s.d.) — lecture modernisée du dossier entier}

\begin{document}

\section*{Résumé}

\begin{resume}
Ce dossier appartient à \emph{Pursuing Stacks} (« À la poursuite des
champs »), le long texte que Grothendieck tape en 1983 et qui part d'une
lettre à Daniel Quillen. L'inventaire de Montpellier le décrit comme le
tapuscrit annoté du chapitre VII, « Linearization of homotopy types »,
accompagné de notes manuscrites. Ce texte a aujourd'hui une édition savante,
celle de Georges Maltsiniotis, et c'est à elle qu'il faut aller pour la
mathématique du dossier : rien de ce qu'elle contient n'est repris ici.

Une seule feuille est lue, parce que l'édition ne la reproduit pas, et elle
ne contient pas de mathématique. Ce sont des consignes tapées à la machine et
signées, pour faire reproduire une tranche du tapuscrit --- les pages 259 à
381, que Grothendieck appelle « Modelizing Story ». Il demande trois copies
et le retour de l'original chez lui ; une copie part pour Bangor, au pays de
Galles, chez Ronnie Brown ; une autre pour Paris, chez Zoghman Mebkhout ; la
troisième est remise à Contou-Carrère. Il signale aussi qu'une feuille a été
tapée recto verso par erreur et qu'il ne faut pas en oublier le verso.

C'est la feuille la plus modeste qui soit, et c'est pour cela qu'elle
instruit. \emph{Pursuing Stacks} n'a pas été publié en son temps : il a
circulé en photocopies, tranche par tranche, à mesure que les pages sortaient
de la machine, vers une poignée de correspondants. Cette feuille montre le
mécanisme à l'œuvre --- qui reçoit, sous quelle forme, avec quelles
précautions pour que rien ne soit coupé ni perdu --- et elle montre que le
texte était lu pendant qu'il s'écrivait. Ronnie Brown, à Bangor, est
justement l'un des correspondants par qui le texte a ensuite été connu.

Pour la suite, les noms à chercher sont ceux du texte lui-même :
\emph{Pursuing Stacks}, ses catégories test et ses modélisateurs, et la
théorie de l'homotopie qu'il annonce.

\keywords{Pursuing Stacks, manuscript circulation, typescript, Ronnie Brown}
\end{resume}

\section*{Ce que couvre cette lecture}
\pagerange{63}{63}

Le dossier compte quatre-vingt-six pages. Une seule, la page~63, est
transcrite, au titre de l'issue~46 du projet : le reste --- le tapuscrit du
chapitre et ses annotations, les notes manuscrites et, parmi les pages 61
à 80, des listes portées au dos de certaines feuilles --- relève de l'édition de \emph{Pursuing Stacks} par
G.~Maltsiniotis, à laquelle cette lecture renvoie sans rien en restituer.

La feuille n'est pas datée.\footnote{La datation « été 1983 » est une
inférence : le post-scriptum place l'envoi peu avant le début d'août, et les
pages 259 à 381 du tapuscrit ne peuvent guère avoir existé avant 1983, année
où le texte est commencé. La feuille ne porte pas d'année ; l'inventaire date
le dossier entier de [1983].} Son post-scriptum la situe quelques jours avant
le début d'un mois d'août, vraisemblablement celui de 1983. Elle concerne les
pages 259 à 381, alors que le dossier, selon l'inventaire, est celui du
chapitre~VII (pages 555 à 593, daté du 22 octobre au 12 novembre) : pourquoi
elle a été rangée là n'est pas dit, et on ne le devine pas
ici.\footnote{Elle a pu servir de chemise ou de brouillon, être réemployée,
ou avoir été classée là après coup ; aucune de ces hypothèses n'est appuyée
par la transcription. La feuille est par ailleurs numérisée à l'envers.}

\section*{Les consignes de reproduction (page 63)}
\pagerange{63}{63}

Voici la feuille, rendue en français courant ; elle est tapée à la machine,
et seules la formule finale et la signature sont de sa main.

\begin{quote}
Recommandations pour la photocopie\footnote{Il a d'abord tapé
« photocopie », puis l'a barré et remplacé par « xérographie », le mot du
temps pour le même procédé.} de « Modelizing Story », pages 259 à 381 :

1) Faire attention au centrage, pour ne pas couper une partie du texte.

2) Faire attention aux pages 274 et 275,\footnote{Un astérisque renvoie en
bas de page : « Plus trois pages de table des matières. » La tranche à
reproduire comprend donc aussi trois pages de table, non numérotées dans le
compte 259--381.} qui ont été tapées par erreur au recto et au verso d'une
même feuille (ne pas oublier, donc, la page 275 au verso !). Faire une copie
supplémentaire de la page 275.

3) Faire la copie en trois exemplaires,\footnote{Plusieurs mots barrés sur
cette ligne et la suivante, dont deux illisibles ; les ajouts « Faire la »,
« l'original » et « copies » les remplacent. Le sens n'en dépend pas.} me
renvoyer l'original à mon domicile (Les Aumettes\footnote{La suite de
l'adresse est omise, comme dans la transcription : c'est une adresse privée.}),
et envoyer deux copies aux adresses suivantes :
\end{quote}

\begin{quote}
Ronnie Brown, School of Mathematics, University College of North Wales,
Bangor, Gwynned\footnote{Sic, pour Gwynedd, le comté gallois.} LL57~2UW,
Angleterre\footnote{Sic : Bangor est au pays de Galles. Le code postal est
tapé « LL 57 2 U W ».}
\end{quote}

et

\begin{quote}
Zoghman Mebkhout, Paris\footnote{Adresse privée omise, comme dans la
transcription ; seule la ville est gardée.}
\end{quote}

\begin{quote}
Remettre la troisième copie à Monsieur Contou-Carrère.

Avec mes remerciements,

A.~Grothendieck

P.-S. Comme je serai absent de chez moi jusque vers le 6 août, il vaudrait
mieux ne renvoyer l'original que vers le 4 ou le 5 août.\footnote{La page a
« aoùt », avec un accent grave sur le \emph{u} --- une faute de frappe,
signalée \emph{sic} par la transcription, qui marque aussi la lecture comme
incertaine.}
\end{quote}

\section*{Ce que la feuille montre}
\pagerange{63}{63}

Le destinataire n'est pas nommé : c'est un service de reproduction, ou une
personne chargée de le solliciter.\footnote{Rien sur la feuille n'identifie ce
destinataire ; l'hypothèse d'un service du département de mathématiques de
Montpellier est naturelle mais n'est pas attestée ici.} Trois choses s'en
dégagent.

D'abord, la diffusion se fait par tranches. La tranche de 123 pages, plus la
table, est envoyée alors que le tapuscrit continue d'avancer : le dossier
lui-même porte un chapitre tapé à l'automne, près de deux cents pages plus
loin. Les
lecteurs suivent donc le texte en cours d'écriture, ce qui convient à un
ouvrage qui se présente comme un journal de recherche plutôt que comme un
traité.

Ensuite, les destinataires. Ronnie Brown, à Bangor, travaillait sur les
groupoïdes et l'homotopie non abélienne, et c'est avec lui que Grothendieck
correspondait sur ces questions.\footnote{Cité de mémoire : la
correspondance Brown--Grothendieck des années 1980 et le rôle de Bangor dans
la circulation des copies de \emph{Pursuing Stacks} sont connus, mais n'ont pas
été vérifiés pour cette lecture.} Zoghman Mebkhout est le mathématicien de la
correspondance de Riemann--Hilbert, dont Grothendieck prendra la défense
quelques années plus tard dans \emph{Récoltes et Semailles}.\footnote{Cité de
mémoire, comme la note précédente.} Contou-Carrère est un collègue de
Montpellier, à qui la copie est « remise » et non envoyée.

Enfin, le soin matériel : le centrage, pour qu'aucune marge ne mange le
texte ; la feuille tapée recto verso par erreur, dont le verso risquait de
disparaître à la photocopie ; une copie supplémentaire de cette page, sans
doute pour la remplacer dans l'original. Rien de cela n'est mathématique,
mais c'est de cette façon qu'un texte qui ne fut imprimé qu'une quarantaine
d'années plus tard a pu être lu dès 1983.

\end{document}
